\section{Order and block structure in detail}\label{sec:structure}
This section records elementary consequences of the host definition
that are useful when interpreting the certificate classification.
The binary word of a node $(d,p)$ consists of the first $d$ bits in
the length-$h$ binary expansion of an appropriate leaf label.
An internal signed node is represented by the corresponding word,
together with its sign.  The leaf at depth $h$ has a single type.

\begin{lemma}[Tree comparison by words]\label{lem:wordorder}
Take two distinct variable blocks of the same axis. If their node
words have an incomparable prefix, their order is determined by the
first bit on which they differ. If one word is a proper prefix of
the other, the ancestor precedes the descendant precisely when its
sign is the opening sign for that axis. The opening signs are minus
for rows and plus for columns. If the two words agree, the blocks
are the two signs at one internal node, and the opening block precedes
the closing block.
\end{lemma}
\begin{proof}
At any internal vertex the recursive row list has the form
\[
 R_d^-(p),\quad \text{left subtree},\quad
 \text{right subtree},\quad R_d^+(p),
\]
whereas the column list exchanges the two signed blocks. If the
words separate into the two different child subtrees, the left
subtree precedes the right regardless of the signs within either.
If one word names an ancestor of the other, the ancestor's opening
block occurs before the descendant's entire subtree and its closing
block after it. Finally, equal words can name distinct blocks only
at an internal node with opposite signs; the displayed recursive
lists place the opening sign before the closing sign. These three
cases exhaust pairs of distinct variable blocks.
\end{proof}

\begin{proposition}[A uniform quotient principle]\label{prop:uniformquotient}
Let $F$ be an $s\times t$ binary matrix with no repeated row
signature and no repeated column signature. Let $M$ be any finite
block-constant ordered binary matrix whose row blocks and column
blocks are nonempty consecutive intervals. Then
\[
 N_F(M)=\sum_{(r,c)\in\mathcal Q_F(M)}
      \prod_{i=1}^{s}|R_{r_i}|\prod_{j=1}^{t}|C_{c_j}|,
\]
where $\mathcal Q_F(M)$ consists of the increasing row-block and
column-block tuples realizing the entries of $F$ on the quotient.
Moreover, $M$ is $F$-free if and only if its quotient has no such
tuple.
\end{proposition}
\begin{proof}
If two chosen physical rows belonged to the same block, then they
would have identical entries on every chosen column, contradicting
the distinctness of the corresponding rows of $F$. The column
argument is symmetric. It follows that each physical copy projects
to a unique quotient tuple with strictly increasing block indices.
Conversely, any tuple of quotient blocks realizing $F$ admits a
physical copy by choosing one position from each participating
block. Every such choice preserves the order, because the blocks
are intervals; it also preserves all entries, because entries are
block-constant. The choices are independent and their number is the
product of the participating block sizes. Summing over the quotient
tuples gives the identity and its freeness consequence.
\end{proof}

There are two independent hypotheses in Proposition~\ref{prop:uniformquotient}.
The consecutive-interval requirement is responsible for preservation
of the ordered embedding, while distinct row and column signatures
ensure that the selected type tuples are strictly increasing.
For patterns with repeated signatures one may have several physical
rows in the same type, and a multiplicity formula requires binomial
factors and within-block order. The simple product formula suffices
for $H_4$ and for the general twin-free patterns considered in
Appendix~\ref{app:smallwitness}.

\subsection{Exact support signatures}
For a mode $t$ and its two marker groups, write
$\rho_1(t),\rho_2(t)$ for row-role predicates and
$\kappa_1(t),\kappa_2(t)$ for column-role predicates.
The marker rows of group $g$ have entry one against a non-marker
column exactly when $\kappa_g(t)$ holds there. Similarly, marker
columns of group $g$ have entry one against a non-marker row
exactly when $\rho_g(t)$ holds.  The corner of a single mode is
$J_2-I_2$, while the corner between different modes is zero.
The parity index in the $W$ modes is therefore a genuine structural
datum. It separates the two candidate child branches using only
one bit, and it is recovered from the parity of the terminal leaf.

\begin{lemma}[Propagation frames realize the prescribed order]
\label{lem:framesorder}
At every depth $1\le i\le h$ and for every node $p_i$ with parent
$p_{i-1}$, the four row and column pairs displayed in
Section~\ref{sec:distance} are increasing tuples on their respective
axes. Their marker restrictions and their incidences with the dummy
blocks agree with the top twelve entries of $H_4$.
\end{lemma}
\begin{proof}
The first two selected positions are the two groups of one marker
mode, both before the variable and dummy positions. The order of
parent and child among variable rows is ancestor--descendant for
minus and descendant--ancestor for plus; among columns it is the
reverse. The first column in a $V$ frame is the column dummy,
which precedes the tree. The second row in a $W$ frame is the row
dummy, which follows it. Consequently the four increasing row and
column tuples are exactly those displayed in the forcing table.
The marker corner is the off-diagonal matrix $J_2-I_2$.
Membership in the two row roles gives the lower-left identity,
and membership in the two column roles gives the upper-right
identity. The dummy incidences are one by the entry definition.
The corresponding body entries may change under edits, yielding
precisely the implications shown in the forcing table.
\end{proof}

\subsection{Where the two signs meet}
At depths below $h$, plus and minus designate different physical
row and column blocks even when they carry the same node label.
Their diagonal entries can therefore consistently have opposite
values. At the leaf level there is only one block $R_h(p)$ and only
one block $C_h(p)$; the selected offsets are both zero. Thus the
positive and negative inductive requirements apply to the same
physical matrix cell, not simply to two equal labels. The
identification is the final step required by the edit-distance
argument. It also explains why the protected selection must use
one common leaf $z=a\oplus b$ on both axes.
